\section{Memory laws uniform in the size of the graph}
\label{sec:v28-growing}
We now combine localized acquisition with a quantitative expansion
condition. All scalar data remain fixed as in
Section~\ref{sec:v28-readout}. Choose first-block charts with
\[
 \Prob(A_e)=\beta(a)\ge\beta_0>0,\qquad
 \Prob\{A_e,X_{e,\ell}\in U\}
            \le h_\ell D_\ell(a)^{-1}\operatorname{vol}_\ell(U).
\]
The constants $\beta_0,h_\ell,L_\ell$ depend on those scalar data,
not on the graph. Their existence follows from
Lemmas~\ref{lem:v25-regular-box} and \ref{lem:v26-first-block} and
the scalar recovery argument. The prior may be any fixed full-support
prior admitted there; it need not have a density.

\subsection{Random retention at a compulsory level}
The random graph formed by retaining the edges for which $A_e$ occurs
has independent Bernoulli edge indicators. For a boundary $j$ and an
integer $r_0\ge1$, put
\begin{equation}\label{eq:v28-retained-cut}
 \begin{split}
 E_{j,r_0}&=\{\,|\delta(S)\cap\{e:A_e\}|\ge r_0
                              \text{ for every }|S|=j\,\},\\
 N_{j,r_0}(G)&=\sum_{|S|=j}\binom{|\delta(S)|}{r_0}.
 \end{split}
\end{equation}
As usual, a binomial coefficient is zero when its lower index exceeds
its upper index. If $\rho=\Prob(E_{j,r_0})>0$, then each such cut has
at least $r_0$ edges and $N_{j,r_0}(G)>0$.

\begin{proposition}[An explicit localized tail]
\label{prop:v28-explicit-tail}
Suppose $\rho>0$. For every positive scalar prefix length $\ell$,
put $k=r_0\ell$ and
\begin{equation}\label{eq:v28-tail-coefficient}
 A_\ell=N_{j,r_0}(G)\,
       \frac{h_\ell^{r_0}}{D_\ell(a)^{r_0}}\,
       \omega_k\bigl(\sqrt{2P}\,L_\ell\bigr)^k.
\end{equation}
Under the enlarged menu, every controller whose decoder has only its
label and the visited-set descriptor satisfies
\begin{equation}\label{eq:v28-explicit-tail}
 \max_i\E\mathcal L_i^\diamond
       \ge\frac{k}{k+2}\rho^{1+2/k}(A_\ell M)^{-2/k}.
\end{equation}
Its scheduler may read its full past. The bound involves neither the
probability that all edges are regular nor a factor $1/(v-1)$.
\end{proposition}
\begin{proof}
At each $S$ use all $r_0$-element subsets of its active cut as witnesses,
assigning prefix length $\ell$ to every selected edge. The event in
\eqref{eq:v28-robust-event} is exactly $E_{j,r_0}$. The selected-edge
density is $h_\ell^{r_0}/D_\ell^{r_0}$, while
Lemma~\ref{lem:v28-local-recovery} bounds the recovery norm by
$\sqrt{2P}L_\ell$. There are $N_{j,r_0}(G)$ witnesses. Apply
Corollary~\ref{cor:v28-power-tail}. All coordinates used in a witness
are actual acquired coordinates on its own selected events; no event
on an edge outside the witness has been inserted into its density.
\end{proof}

\begin{lemma}[Retention of all balanced cuts]\label{lem:v28-retention}
Suppose every set of size $j=\lfloor v/2\rfloor$ has at least $h v$
crossing edges. For $r_0=\lfloor\beta_0 h v/2\rfloor\ge1$,
\begin{equation}\label{eq:v28-retention-probability}
 \Prob(E_{j,r_0})\ge
       1-\binom vj\exp(-\beta_0 h v/8)
       \ge1-\exp\{-(\beta_0h/8-\log 2)v\}.
\end{equation}
Here and below $\log$ denotes the natural logarithm when its base is
not indicated.
\end{lemma}
\begin{proof}
For a fixed cut let $X$ count its retained edges and let
$\mu=\E X\ge\beta_0 h v$. Independence gives
$\E 2^{-X}=\prod_e(1-p_e/2)\le e^{-\mu/2}$, where $p_e$ are the
retention probabilities on this cut. Markov's inequality yields
\[
 \Prob\{X\le\mu/2\}
       \le \exp\{-(1-\log 2)\mu/2\}\le e^{-\mu/8}.
\]
The last inequality uses $\log 2<3/4$. Since $r_0\le\mu/2$,
the same estimate bounds $\Prob\{X<r_0\}$. Sum over all sets at
level $j$ and use $\binom vj\le2^v$. The retained-edge law is fixed
before an order is chosen. Thus this simultaneous estimate applies
even to a visited set selected using all the revealed reports.
\end{proof}

\subsection{A uniform adaptive bit law}
\begin{theorem}[Size-uniform collision-dependent memory]
\label{thm:v28-size-uniform}
Fix numbers $d,h>0$ with
\begin{equation}\label{eq:v28-expansion-margin}
                       \beta_0h/8>\log 2.
\end{equation}
Consider all loopless multigraphs with $v$ vertices, $P\le d v$ edges,
and $|\delta(S)|\ge h v$ whenever $|S|=\lfloor v/2\rfloor$.
There are constants $v_0,C_-,C_+>0$, depending only on these numbers
and the fixed scalar data, such that for every such graph with
$v\ge v_0$, every admitted calibration and $0<\varepsilon\le1$,
\begin{equation}\label{eq:v28-additive-bit-law}
 \begin{split}
 r_0 b(\varepsilon,a)-C_-v
 &\le B^{\diamond,\rm av}_G(\varepsilon,a)\\
 &\le B^{\diamond,\rm max}_G(\varepsilon,a)
 \le \operatorname{cw}(G)\{b(\varepsilon,a)+C_+\},
 \qquad r_0=\lfloor\beta_0 h v/2\rfloor .
 \end{split}
\end{equation}
The lower bound continues to hold for a full-history scheduler when
the decoder receives only its label and the visited-set descriptor.
In particular, there are $\varepsilon_0,c,C>0$, independent of the
graph and calibration, such that
\begin{equation}\label{eq:v28-multiplicative-bit-law}
 c v b(\varepsilon,a)
 \le B^{\diamond,\rm av}_G(\varepsilon,a)
 \le B^{\diamond,\rm max}_G(\varepsilon,a)
 \le C v b(\varepsilon,a)
                   \qquad(0<\varepsilon\le\varepsilon_0).
\end{equation}
All upper bounds use an admissible deterministic order and one causal
joint-label filter. Exact additive exponent collisions are included.
\end{theorem}
\begin{proof}
Put $\gamma=\beta_0h/4$ and
$c_0=\beta_0h/8-\log 2>0$. Choose $v_0$ so that
$r_0\ge\gamma v\ge1$ and $e^{-c_0v}\le1/4$ for $v\ge v_0$.
Lemma~\ref{lem:v28-retention} then gives $\rho\ge3/4$.
Proposition~\ref{prop:v28-uniform-upper} proves the upper bound in
\eqref{eq:v28-additive-bit-law}; the middle inequality follows from
the common independent tape law.

For the lower bound fix a positive prefix length $\ell$, and suppose
first that one controller with $M=2^B$ has maximum expected boundary
loss at most $\varepsilon$. At the fixed balanced boundary, Markov's
inequality gives
\[
 \Prob\{E_{j,r_0},\mathcal L_j^\diamond\le4\varepsilon\}
                         \ge\rho-1/4\ge1/2.
\]
The small-ball proof of Proposition~\ref{prop:v28-explicit-tail}
therefore implies $1/2\le M A_\ell(4\varepsilon)^{k/2}$.
We record the dimension-dependent ball volume rather than replacing
it by an exponential bound with a hidden dimension constant. The
Gaussian integral gives, for every $k\ge1$,
\begin{equation}\label{eq:v28-ball-volume}
                         \omega_k\le(2\pi e/k)^{k/2}.
\end{equation}
Indeed, integration over the ball of radius $\sqrt k$ is at least
$e^{-k/2}\omega_k k^{k/2}$ and at most
$\int_{\R^k}e^{-|x|^2/2}\,dx=(2\pi)^{k/2}$.
Substituting \eqref{eq:v28-tail-coefficient}, taking base-two
logarithms, and using $k=r_0\ell$ yields
\begin{equation}\label{eq:v28-finite-size-bits}
 \begin{split}
 B\ge{}&r_0\left\{\log_2D_\ell(a)
                      +\frac\ell2\log_2(1/\varepsilon)\right\}
                  -\log_2\{2N_{j,r_0}(G)\}\\
       &-r_0\log_2h_\ell
         -\frac{r_0\ell}{2}
              \log_2\left(\frac{16\pi e L_\ell^2P}{r_0\ell}\right).
 \end{split}
\end{equation}
The factor $16$ includes both the probability $1/2$ of the local menu
and the threshold $4\varepsilon$ in Markov's inequality.

Now $N_{j,r_0}(G)\le2^{v+P}$, $P/r_0\le d/\gamma$ and
$r_0\le P\le dv$. Set
\[
 h_* =\max(1,h_1,\ldots,h_p),\qquad
 L_* =\max(1,L_1,\ldots,L_p),
\]
and choose
\[
 K_* =\log_2(h_*p!)+\frac p2
                 \log_2\max\{1,16\pi e L_*^2d/\gamma\}.
\]
Using $V_\ell\le\ell!D_\ell$ in
\eqref{eq:v28-finite-size-bits}, then maximizing over the positive
indices in \eqref{eq:v28-single-bit-profile}, gives
\[
             B\ge r_0b(\varepsilon,a)
                            -\{2+d+dK_*\}v.
\]
If the maximum is at the zero index, this follows already from
$B\ge0$. An unattained risk infimum causes no difficulty: apply the
bound at $\varepsilon+\delta$ to approximating controllers and let
$\delta\downarrow0$. The profile is a finite continuous maximum as
a function of positive $\varepsilon$. This proves the lower bound
with $C_-=2+d+dK_*$.

Finally $r_0\ge\gamma v$, $\operatorname{cw}(G)\le dv$, and
$b(\varepsilon,a)\ge\tfrac12\log_2(1/\varepsilon)$. Choose
$\varepsilon_0$ so that this last lower bound is at least
$\max\{1,2C_-/\gamma\}$. The additive lower loss is then absorbed
into half the leading term, and the upper additive term is absorbed
using $b\ge1$. This proves \eqref{eq:v28-multiplicative-bit-law},
for example with $c=\gamma/2$ and $C=d(1+C_+)$. None of these
constants depends on $v$, calibration or the memory budget.
\end{proof}

This theorem is not an exact-constant cutwidth formula. It supplies
uniform upper and lower memory orders over a family of increasing
experiments; the additive form \eqref{eq:v28-additive-bit-law} gives
more information at moderate resolutions. It is also not a uniform
statement for the tensor-only loss. That loss remains governed, for
each fixed graph, by the earlier theorems. The present enlarged menu
retains the tensor task with fixed positive weight and explicitly
adds locally readable tests from the same future test space.

\subsection{Existence of linear-size families}
The hypotheses describe nonempty families of arbitrary size; this
need not be assumed from a graph construction theorem.

\begin{proposition}[Linear-size graphs with the required balanced cuts]
\label{prop:v28-graph-existence}
For every integer $D>32\log 2$ and every integer $s\ge1$ there is a
loopless bipartite multigraph on $v=2s$ vertices with $P=Ds$ edges
such that every $s$-vertex set has at least $Ds/4=Dv/8$ crossing
edges. Consequently, for every admitted scalar experiment one can
choose a fixed integer $D$ with $\beta_0D/64>\log 2$ and obtain
arbitrarily large graphs satisfying
Theorem~\ref{thm:v28-size-uniform}, with $d=D/2$ and $h=D/8$.
\end{proposition}
\begin{proof}
Take two parts of size $s$. From every left vertex choose $D$ right
endpoints independently and uniformly, allowing repetitions and
retaining the resulting edges as distinct labelled edges. Fix an
$s$-vertex set containing $a$ left vertices and $s-a$ right vertices.
Each outgoing edge of a selected left vertex crosses with probability
$a/s$, and each outgoing edge of an unselected left vertex crosses
with probability $(s-a)/s$. These indicators are independent. Their
sum has mean
\[
          \mu=\frac D s\{a^2+(s-a)^2\}\ge Ds/2.
\]
The exponential estimate in the proof of
Lemma~\ref{lem:v28-retention} shows that this cut has fewer than
$Ds/4$ edges with probability at most $e^{-Ds/16}$. There are at most
$4^s$ candidate sets. Hence the probability of any failed balanced cut
is at most $\exp\{(2\log 2-D/16)s\}<1$. Some realization has the
asserted property. There are exactly $Ds$ edges and no loops. In the
last assertion choose $D$ satisfying both displayed strict inequalities;
this is possible because $\beta_0>0$. The graphs have bounded average
degree for this fixed $D$. A uniform bound on their maximum degree is
not claimed or needed.
\end{proof}

\begin{corollary}[Simultaneous graph growth and calibration collision]
\label{cor:v28-joint-limit}
Let $a(\theta)$ be an admitted analytic calibration path and suppose
$V_\ell(a(\theta))\asymp\theta^{\eta_\ell}$ for each nonzero
profile term. Put $\eta_0=0$, omit identically zero terms, and set
$U=\log_2(1/\theta)$ and
$w(s)=\max_{0\le\ell\le p}(\ell s-\eta_\ell)$.
For graphs in Theorem~\ref{thm:v28-size-uniform} and a fixed compact
interval $I\subset(0,\infty)$, uniformly for $s\in I$,
\begin{equation}\label{eq:v28-joint-limit}
 \gamma v U w(s)-C_1v
 \le B_G^{\diamond,\rm av}(\theta^{2s},a(\theta))
 \le B_G^{\diamond,\rm max}(\theta^{2s},a(\theta))
 \le \operatorname{cw}(G)U w(s)+C_2v.
\end{equation}
The constants are independent of the graph and $\theta$ sufficiently
small. In particular, graph size may tend to infinity at an arbitrary
rate as the calibration approaches the collision locus.
\end{corollary}
\begin{proof}
The scalar contact-order analysis gives
$b(\theta^{2s},a(\theta))=Uw(s)+O(1)$ uniformly on $I$; there are
only finitely many capped determinant terms. Substitute this expression
in \eqref{eq:v28-additive-bit-law}, using
$r_0\ge\gamma v$, $r_0\le dv$, $w(s)\ge0$ and
$\operatorname{cw}(G)\le dv$. Every bounded scalar remainder is
multiplied by at most a constant times $v$, which gives exactly the
remainders stated in \eqref{eq:v28-joint-limit}. No graph-dependent
constant is exchanged with the collision limit.
\end{proof}
